% Generated by opera/run/render_borrowing_pairs_review.py from
% opera/configs/borrowing_hypotheses.yaml; do not edit by hand.
\begingroup\scriptsize\setlength{\tabcolsep}{3pt}
\begin{longtable}{@{}p{0.03\linewidth}p{0.15\linewidth}p{0.15\linewidth}p{0.13\linewidth}p{0.07\linewidth}p{0.41\linewidth}@{}}
\caption{\textbf{Prespecified borrowing pairs and negative controls.} Fixed before the
test set was opened and reviewed by the study haematologists. Source: the disease
switched on alone; target: the disease predicted; outcomes: the clinical domain
tested. Rating: the study haematologists' expected strength, given blind to the
rationale and to whether the pair was proposed as a hypothesis or a control. Pairs were labelled against a practical margin of 0.25 AUC points.}\label{tab:borrowing-pairs}\\
\toprule & Source & Target & Outcomes & Rating & Rationale (evidence)\\\midrule\endfirsthead
\toprule & Source & Target & Outcomes & Rating & Rationale (evidence)\\\midrule\endhead
\bottomrule\endfoot
\multicolumn{6}{@{}l}{\textit{Expected to share signal}}\\
H1 & Chronic lymphocytic leukaemia (CLL) & Small lymphocytic lymphoma (SLL) & all outcomes & -- & One disease in two presentations (blood vs lymph nodes), managed by the same criteria and treatments. (WHO 5th edition classification (Alaggio et al., Leukemia 2022))\\
H2 & Extranodal marginal-zone lymphoma (EMZL) & Splenic marginal-zone lymphoma (SMZL) & disease course and later malignancy & -- & The marginal-zone lymphomas share prognostic factors: the MZL-IPI (LDH, haemoglobin, lymphocyte and platelet counts, subtype) was built for and validated across all subtypes, and advanced disease is treated with the same rituximab-based regimens. First-line treatment often differs (radiotherapy for localised EMZL; rituximab alone or splenectomy for SMZL). (Arcaini et al., EClinicalMedicine 2024, doi:10.1016/j.eclinm.2024.102592)\\
H3 & Multiple myeloma (MM) & Chronic lymphocytic leukaemia (CLL) & infection and antimicrobial use & -- & Both cause secondary humoral immune deficiency (hypogammaglobulinaemia in up to 85\% of CLL during the disease course, and intrinsic or treatment-related in myeloma), with the same immunoglobulin-replacement indications. (Frontiers in Immunology 2014, doi:10.3389/fimmu.2014.00626; Frontiers in Immunology 2016, doi:10.3389/fimmu.2016.00317)\\
H4 & Chronic lymphocytic leukaemia (CLL) & Multiple myeloma (MM) & infection and antimicrobial use & -- & The same, in the other direction. (As above)\\
H5 & Burkitt lymphoma (BL) & Lymphoblastic lymphoma (LBL) & neutrophil, white-cell and platelet findings & -- & Both receive short, highly intensive regimens (CODOX-M/IVAC, hyper-CVAD, ALL-type protocols) with marked myelosuppression despite growth-factor support. Holds whether or not the LBL cohort includes adult ALL (open question for the haematologists). (Reviews of adult Burkitt lymphoma treatment (hyper-CVAD, CODOX-M/IVAC toxicity))\\
H6 & Burkitt lymphoma (BL) & Lymphoblastic lymphoma (LBL) & metabolic and cell-turnover markers & -- & Rapid tumour turnover and intensive induction put both at risk of tumour lysis (potassium, phosphate, urate, creatinine). (Reviews of adult Burkitt lymphoma treatment (tumour lysis precautions))\\
H7 & Peripheral T-cell lymphoma, NOS (PTCL) & T-cell lymphoma, unspecified (TCL) & disease course and later malignancy & -- & Unspecified T-cell lymphoma is managed like peripheral T-cell lymphoma, not otherwise specified (CHOP-based therapy), with a similarly poor course. (WHO 5th edition classification (Alaggio et al., Leukemia 2022))\\
H8 & Peripheral T-cell lymphoma, NOS (PTCL) & Angioimmunoblastic T-cell lymphoma (AITL) & disease course and later malignancy & -- & Angioimmunoblastic T-cell lymphoma (nodal T-follicular-helper lymphoma in WHO 5) is treated with the same CHOP-based first-line regimens as PTCL, not otherwise specified, with a similarly poor course; the two are often hard to tell apart histologically. (WHO 5th edition classification (Alaggio et al., Leukemia 2022))\\
H9 & Chronic lymphocytic leukaemia (CLL) & Lymphoplasmacytic lymphoma (LPL) & cardiovascular and thrombotic events & -- & BTK inhibitors, used in CLL and in Waldenstrom macroglobulinaemia, cause atrial fibrillation and bleeding. But first-line Waldenstrom treatment in Denmark has mostly been chemoimmunotherapy (e.g. bendamustine-rituximab), so a null result here would be unsurprising. (Brown et al., Haematologica 2017 (pooled AF in ibrutinib trials), PMC5622864)\\
\multicolumn{6}{@{}l}{\textit{Negative controls}}\\
C1 & Multiple myeloma (MM) & Hodgkin lymphoma (HL) & all outcomes & -- & No shared biology, treatment or complication profile (young patients on ABVD-type therapy vs older patients on myeloma regimens); MM is the second-largest source, so this also tests whether a large source helps regardless of relevance. (--)\\
C2 & Hairy-cell leukaemia (HCL) & Diffuse large B-cell lymphoma (DLBCL) & disease course and later malignancy & -- & No shared biology or course (indolent, cladribine-treated vs aggressive, R-CHOP-treated). (--)\\
C3 & Anaplastic large-cell lymphoma (ALCL) & Chronic lymphocytic leukaemia (CLL) & all outcomes & -- & No shared biology, treatment or complication profile (a T-cell lymphoma treated with CHOP-type therapy vs an indolent B-cell leukaemia). (--)\\
\end{longtable}
\endgroup
